\chapter*{Abstract}

How do social norms emerge in multi-agent systems where Large Language Models (LLMs) interact as autonomous agents? As AI agents proliferate in collaborative and social settings—from scientific co-authoring to policy deliberation—understanding norm emergence in these architectures has become critical for HCI and social computing. We present a computational study of norm emergence in populations of LLM-based agents across three interaction regimes: deliberation, negotiation, and coordination. Our findings reveal that coordination tasks produce the strongest norm convergence (73\% behavioral consistency), while deliberation tasks foster diverse epistemic norms with significant within-population variance. We identify four distinct emergence trajectories and demonstrate that agents exhibiting higher theory-of-mind reasoning are more likely to initiate norm-breaking behavior, which can cascade into population-wide shifts. We further show that agent identity (persona consistency) and prompt framing significantly moderate norm emergence speed and stability. This work contributes: (1) empirical evidence of norm emergence patterns in LLM multi-agent systems, (2) a typology of emergence trajectories with associated conditions, and (3) design implications for HCI systems deploying multi-agent AI. Our findings suggest that norm governance in multi-agent AI systems requires explicit structural interventions rather than relying on emergent self-organization.

\textbf{Keywords:} norm emergence, multi-agent systems, Large Language Models, social norms, collective behavior, HCI, AI agents